% =========================================================================
% Appendix B: Prompt Engineering for LLM-based Visualization Analysis (v2)
% Appendix C: Topic Selection (Follow-on Study)
% Two-column IEEE VIS / CHI style.
% Requires in the preamble:
%   \usepackage{booktabs}
%   \usepackage{array}
%   \usepackage{listings}   % for prompt-listing blocks in \S app:prompts
%   Recommended lstset (add once in preamble):
%     \lstset{basicstyle=\ttfamily\footnotesize, breaklines=true,
%             frame=single, columns=fullflexible, keepspaces=true,
%             aboveskip=4pt, belowskip=4pt}
% Uses \cite{lin1989ccc}.
% Tables use tabular* with \extracolsep{\fill} so cells stay content-tight
% but inter-column gaps expand to fill the full column / page width.
% Narrow tables: \begin{table}  + tabular*{\linewidth}  (single column).
% Wide tables:   \begin{table*} + tabular*{\textwidth} (spans both columns).
%
% v2 changes (Apr 2026):
%   - Unified E1+E2 into a single ablation; former V3 is renamed V1 and
%     enters as the "extended" variant on top of V0+TWCA.
%   - Dropped V0+CA and V0+TCA (redundant or superseded cells) from
%     the variant list.
%   - Added V0+TW, V0+TWC, V0+TWA cells (TW-anchored triples).
%   - Rationale: with three follow-on analyses planned (VC prediction,
%     topic selection, scene-graph extraction), the appendix title is
%     generalized from "Visual Complexity Scoring" to "LLM-based
%     Visualization Analysis". The VC-prediction study below is the
%     first of the three.
% =========================================================================

\appendix
\section{Prompt Engineering for LLM-based Visualization Analysis}
\label{app:vc-prompts}

We investigate how the structure of a large multimodal LLM prompt affects
the model's ability to reproduce human-rated \textbf{Visual Complexity (VC)}
scores for static data visualizations. Starting from a minimal baseline
prompt (\textbf{V0}), we incrementally layer four modular components ---
\textbf{T}opics, \textbf{C}alibration, \textbf{A}nchors, and
\textbf{W}eighted scoring --- and report a single unified ablation over ten
variants. An eleventh variant, \textbf{V1}, extends the most heavily
instrumented cell (V0+TWCA) with two additional elements --- per-dimension
anchor supervision and per-dimension numeric output --- and is included as
the upper envelope of the ablation.

\subsection{Ground Truth and Evaluation Set}
\label{app:vc-gt}

The 46-image ground-truth set was sub-sampled from a pool of approximately
700 paired images collected in a prior pair-comparison study. Pairs were
ranked by the absolute difference $|\Delta \mathrm{VC}|$ between their
aggregated human ratings, then binned across the full range of
$|\Delta \mathrm{VC}|$. From each bin we selected roughly equal numbers of
pairs, yielding \textbf{23 pairs (46 images)} that span the full difficulty
spectrum --- from near-ties (hard discrimination) to large-gap pairs (easy
discrimination). Each image carries a normalized human VC rating
$y \in [0, 1]$ derived from the pair study.

Three of the 46 images are reserved as few-shot anchor exemplars
($\mathrm{VC} = 0.22, 0.54, 0.95$) and are presented in-context to every
anchored variant (all cells containing ``A'', plus V1). To prevent leakage,
these three images are \emph{excluded} from all reported metrics, leaving
$n = 43$ test images.

Table~\ref{tab:vistype-dist} shows the vis-type breakdown of the 46 images,
together with the three anchor exemplars.

\begin{table}[h]
  \centering
  \caption{Vis-type distribution of the 46-image pilot set. Three images
    serve as few-shot anchors (marked with $\dagger$).}
  \label{tab:vistype-dist}
  \small
  \begin{tabular*}{\linewidth}{@{\extracolsep{\fill}}lcc@{}}
    \toprule
    Vis Type & Count & Anchor \\
    \midrule
    Bar             & 8 & VisC.503.6$^\dagger$ \\
    Area            & 7 & --- \\
    Glyph           & 7 & --- \\
    Grid            & 6 & InfoVisJ.1149.6(1)$^\dagger$ \\
    Node-link       & 6 & --- \\
    Cont.-ColorPatn & 4 & --- \\
    Point           & 3 & InfoVisJ.619.17$^\dagger$ \\
    Text            & 3 & --- \\
    Line            & 2 & --- \\
    \midrule
    \textbf{Total}  & \textbf{46} & \\
    \bottomrule
  \end{tabular*}
\end{table}

The set spans nine vis types and covers a broad range of visual forms,
from common statistical charts (Bar, Area, Line) to more complex
encodings (Glyph, Node-link, Cont.-ColorPatn). Because the sampling
was driven by $|\Delta \mathrm{VC}|$ coverage rather than vis-type
stratification, counts are unequal; the three least-represented types
(Point, Text, Line, $n \le 3$) are too small for reliable per-type
breakdowns. All results in subsequent sections are therefore reported
over the full test set.

\subsection{Evaluation Metrics}
\label{app:vc-metrics}

We report five agreement metrics, with Lin's Concordance Correlation
Coefficient (CCC)~\cite{lin1989ccc} as the headline measure:
\begin{equation}
\mathrm{CCC}(y, \hat{y}) =
\frac{2\,\rho\,\sigma_y\sigma_{\hat{y}}}
     {\sigma_y^2 + \sigma_{\hat{y}}^2 + (\mu_y - \mu_{\hat{y}})^2},
\end{equation}
where $\rho$ is the Pearson correlation between human ratings $y$ and model
predictions $\hat{y}$. CCC measures agreement with the identity line
$\hat{y} = y$, jointly penalizing (i) low correlation, (ii) mean bias, and
(iii) scale mismatch. It is bounded in $[-1, 1]$ and is the standard metric
for continuous method-comparison studies. We additionally report Pearson
$r$ with its two-sided $p$-value (\texttt{scipy.stats.pearsonr}), Spearman
$\rho$, the coefficient of determination $R^2$ against the identity line
(\texttt{sklearn.metrics.r2\_score},
$R^2 = 1 - \sum_i (y_i - \hat{y}_i)^2 / \sum_i (y_i - \bar{y})^2$, which
can be negative when the predictor underperforms the mean of $y$), MAE,
RMSE, and the signed bias $\overline{\hat{y} - y}$.

\subsection{Prompt Components}
\label{app:vc-components}

Starting from the minimal \textbf{V0} baseline --- VC definition plus the
0--1 output range, no dimensions, no anchors --- we incrementally layer
four modular components:

\begin{itemize}
  \item \textbf{T (Topics / Dimensions):} Enumerates seven complexity
    dimensions (data density, visual encoding, semantics/text, schema,
    color/symbol, aesthetics, cognitive load) as mental checkpoints.
    Output remains a single holistic \texttt{vc\_score}.
  \item \textbf{C (Calibration):} Provides scale-anchoring guidance with
    example chart types mapped to target ranges (plain bar chart
    $\approx 0.25$--$0.40$; dense multi-panel $\approx 0.85$--$0.95$) and
    instructs the model to use the full range.
  \item \textbf{A (Anchors):} Prepends the three labeled anchor images as
    few-shot exemplars before the target. Each anchor carries only its
    final VC score (no per-dimension values).
  \item \textbf{W (Weighted):} Qualitatively re-weights the seven
    dimensions: high weight on density / encoding / schema / cognitive,
    medium on color / aesthetic, low on text.
\end{itemize}

\subsection{Variant Selection}
\label{app:vc-variants}

Rather than run all $2^3$ combinations of $(T, C, A)$, we prune the lattice
on the observation that \textbf{among the single-component cells (V0+T,
V0+C, V0+A), V0+T achieves the highest CCC}, motivating a T-anchored path
through the design space. The final variant list is:

\begin{table}[t]
  \centering
  \caption{Variant lattice for the prompt-engineering ablation.
    Starting from a minimal baseline (\textbf{V0}: VC definition and
    0--1 output only), we layer four modular components:
    \textbf{T}opics (7 complexity dimensions as checkpoints),
    \textbf{C}alibration (example chart types mapped to target ranges),
    \textbf{A}nchors (three few-shot labeled images), and
    \textbf{W}eighted scoring (high/medium/low dimension weights).
    Rather than all $2^3$ combinations, we prune to a T-anchored path:
    V0+T is the strongest single-component cell, so doubles pair T with
    each remaining component; triples build on the winning double
    (V0+TW); V0+TWCA is the full four-component cell.
    \textbf{V1} is an extended variant that additionally supplies
    per-dimension anchor scores and requires per-dimension numeric
    output, beyond what any E1 cell carries.}
  \label{tab:variant-lattice}
  \small
  \begin{tabular*}{\linewidth}{@{\extracolsep{\fill}}ll@{}}
    \toprule
    Level & Variants\\
    \midrule
    Baseline          & V0\\
    Single            & V0+T, V0+C, V0+A\\
    Double (T-only)   & V0+TC, V0+TA, V0+TW\\
    Triple (TW-based) & V0+TWC, V0+TWA\\
    Quadruple         & V0+TWCA\\
    Extended          & \textbf{V1} (V0+TWCA + per-dim anchors \& output)\\
    \bottomrule
  \end{tabular*}
\end{table}

Since V0+T is the strongest single-component cell, the double-feature
sweep pairs T with each of the three remaining components (C, A, W) ---
giving V0+TC, V0+TA, and V0+TW --- rather than the redundant non-T cell
V0+CA. The triple layer is then built on the strongest double (V0+TW) by
adding C or A, yielding V0+TWC and V0+TWA. V0+TWCA is the full
four-component (quadruple) cell.

\textbf{V1} extends V0+TWCA with two additional elements: (a) the three
anchor images now carry full per-dimension reference scores (not just the
overall VC), and (b) the model is required to emit numeric scores for each
of the seven dimensions in addition to \texttt{vc\_score}.

Dropped relative to the original design: V0+CA (redundant two-feature cell
without T) and V0+TCA (superseded by V0+TWC once W emerges as the
strongest pairing with T).

\subsection{Test Conditions}
\label{app:vc-conditions}

\textbf{E1 variants (V0 through V0+TWCA).} Each variant is evaluated under
four independent runs on \texttt{claude-opus-4-6} (primary model):
\texttt{r1} and \texttt{r2} are deterministic ($t = 0$) replicates;
\texttt{r3} and \texttt{r4} use Claude's adaptive extended-thinking mode.
A subset of \texttt{r3} (variants V0+TC, V0+TA) used
\texttt{claude-sonnet-4-6}; these were re-run in \texttt{r4} on
\texttt{claude-opus-4-6} for cross-variant parity. We define
\emph{det\_avg} as the mean of \texttt{r1} and \texttt{r2}, and
\emph{think\_avg} as the mean of \texttt{r3} and \texttt{r4}.

\textbf{V1.} Run twice on \texttt{claude-opus-4-6} at the Anthropic API
default sampling temperature ($t = 1$, not explicitly set) with no
extended thinking and \texttt{max\_tokens = 1500} to accommodate the
longer per-dimension JSON output. We report both replicates and their
mean; V1 per-run values carry more sampling noise than the E1
\emph{det\_avg} values since $t > 0$.

\subsection{Results --- CCC per run}
\label{app:vc-ccc}

Table~\ref{tab:ccc} summarizes CCC across all runs and aggregated
columns; we then walk through the patterns visible in the table.

\paragraph{Single-component gains are small and ordered $T > C > A$.}
Among the three single-component cells, V0+T leads ($0.869$),
V0+C is mid ($0.854$), and V0+A is lowest ($0.852$) --- only
$0.011$ above the V0 baseline ($0.858$). This ordering is what
motivated the T-anchored pruning of the two- and three-feature cells
in Sec.~\ref{app:vc-variants}: pairing topics with each remaining
component is more informative than re-running the lattice without T.

\paragraph{Layering along the T-path is near-flat at $t = 0$.}
V0+T, V0+TC, V0+TA, V0+TW, V0+TWC, V0+TWA, and V0+TWCA all sit
within $0.018$ CCC of each other ($0.854$--$0.872$). Among the
double-feature T-pairs, \textbf{V0+TW is the strongest ($0.872$)},
narrowly above V0+TA ($0.860$) and V0+TC ($0.854$). Adding C or A
on top of V0+TW does \emph{not} help: the two triples lose
$\sim\!0.015$ CCC (V0+TWC $0.856$, V0+TWA $0.857$) and the
quadruple V0+TWCA ($0.869$) does not recover the V0+TW level.
Within the standard ablation lattice, V0+TW is therefore the best
deterministic cell of the entire sweep.

% ---------- Narrow CCC table ----------------------------------------------
\begin{table}[t]
  \centering
  \caption{Concordance Correlation Coefficient (CCC, Lin~1989) per run
    on $n = 43$ non-anchor images, \texttt{claude-opus-4-6}. CCC
    measures agreement with the identity line $\hat{y} = y$, jointly
    penalizing low correlation, mean bias, and scale mismatch; higher
    is better ($\in [-1, 1]$). \texttt{r1}/\texttt{r2} are
    deterministic ($t = 0$) replicates; \texttt{r3}/\texttt{r4} use
    Claude's adaptive extended-thinking mode.
    \emph{det\_avg} averages \texttt{r1}+\texttt{r2};
    \emph{think\_avg} averages \texttt{r3}+\texttt{r4}. V1 is run at
    $t = 1$ (API default, no thinking); its \emph{det\_avg} column
    reports the mean of the two V1 replicates for column parity.
    Dashes indicate runs not conducted for that cell.}
  \label{tab:ccc}
  \small
  \begin{tabular*}{\linewidth}{@{\extracolsep{\fill}}lcccccc@{}}
    \toprule
    Variant & r1 & r2 & \textbf{det\_avg} & r3 & r4 & \textbf{think\_avg}\\
    \midrule
    V0          & 0.858 & 0.859 & \textbf{0.858} & 0.868 & 0.881 & \textbf{0.875}\\
    V0+T        & 0.871 & 0.867 & \textbf{0.869} & 0.857 & ---   & \textbf{0.857}\\
    V0+C        & 0.851 & 0.856 & \textbf{0.854} & ---   & ---   & ---\\
    V0+A        & 0.852 & 0.853 & \textbf{0.852} & ---   & ---   & ---\\
    V0+TC       & 0.854 & ---   & \textbf{0.854} & 0.737 & 0.866 & \textbf{0.802}\\
    V0+TA       & 0.860 & ---   & \textbf{0.860} & 0.849 & 0.834 & \textbf{0.841}\\
    V0+TW       & 0.872 & 0.872 & \textbf{0.872} & ---   & ---   & ---\\
    V0+TWC      & 0.855 & 0.856 & \textbf{0.856} & ---   & ---   & ---\\
    V0+TWA      & 0.860 & 0.855 & \textbf{0.857} & ---   & ---   & ---\\
    V0+TWCA     & 0.868 & 0.871 & \textbf{0.869} & 0.811 & 0.815 & \textbf{0.813}\\
    \textbf{V1} & 0.882 & 0.880 & \textbf{0.881} & ---   & ---   & ---\\
    \bottomrule
  \end{tabular*}
\end{table}

\paragraph{Extended thinking is a double-edged sword.}
The \emph{think\_avg} column tells a non-monotone story: extended
thinking lifts the minimal V0 prompt ($0.858 \to 0.875$) but sharply
hurts heavily-instructed prompts (V0+TWCA: $0.869 \to 0.813$;
V0+TC: $0.854 \to 0.802$). The qualitative weight guidance and the
calibration ranges appear to interact poorly with longer
deliberation. The low V0+TC \texttt{r3} cell ($0.737$) is also
partly a model-capacity confound: that single run used
\texttt{claude-sonnet-4-6} and is recovered by its \texttt{r4} opus
replicate ($0.866$), so the dip reflects model size rather than a
prompt defect.

\paragraph{V1 breaks the $0.87$ ceiling.}
The extended variant V1 (full per-dimension anchor supervision and
required per-dimension numeric output, on top of V0+TWCA) reaches
CCC $= 0.881$ --- the only cell above the $\sim\!0.87$ band that the
standard ablation never crosses. The two V1 replicates agree closely
($0.882$ and $0.880$), so the $t = 1$ sampling noise is small
relative to the gap to V0+TW. We expand on this gain in
Sec.~\ref{app:vc-full} alongside the full agreement metrics.

\subsection{Results --- Full metrics}
\label{app:vc-full}

Table~\ref{tab:full} reports the full agreement metrics for all
variants. The CCC ordering of Sec.~\ref{app:vc-ccc} is broadly
preserved across $r$, $\rho$, $R^2$, MAE, and RMSE, with one
clear-cut conclusion that the multi-metric view sharpens.

% ---------- Wide full-metrics table ---------------------------------------
\begin{table*}[t]
  \centering
  \caption{Full agreement metrics for all variants on $n = 43$
    non-anchor images, \texttt{claude-opus-4-6}. Ground truth is a
    normalized human VC rating $y \in [0,1]$ from a prior
    pair-comparison study. E1 rows (V0 through V0+TWCA) report
    \emph{det\_avg} (mean of two $t = 0$ replicates); the V1 row
    reports the mean of two $t = 1$ replicates (per-run statistics
    averaged). CCC is Lin's Concordance Correlation Coefficient
    (headline metric); $r$ is Pearson; $\rho$ is Spearman; $R^2$ is
    \texttt{sklearn.metrics.r2\_score} against the identity line
    (can be negative if the predictor underperforms the mean of $y$);
    MAE and RMSE are in the raw $[0,1]$ VC scale; Bias is the signed
    mean error $\overline{\hat{y} - y}$. All Pearson $p < 10^{-12}$.
    \textbf{V1} (bold) is the only variant that simultaneously
    maximizes CCC/$r$/$\rho$/$R^2$ and minimizes MAE/RMSE.}
  \label{tab:full}
  \small
  \begin{tabular*}{\textwidth}{@{\extracolsep{\fill}}lcccccccc@{}}
    \toprule
    Variant & $n$ & CCC & $r$ & $\rho$ (Spearman) & $R^2$ & MAE & RMSE & Bias\\
    \midrule
    V0          & 43 & 0.858          & 0.868          & 0.868          & 0.671          & 0.080          & 0.098          & $+0.010$\\
    V0+T        & 43 & 0.869          & 0.879          & 0.872          & 0.692          & 0.078          & 0.095          & $+0.003$\\
    V0+C        & 43 & 0.854          & 0.871          & 0.873          & 0.679          & 0.081          & 0.097          & $+0.033$\\
    V0+A        & 43 & 0.852          & 0.869          & 0.866          & 0.638          & 0.083          & 0.103          & $-0.012$\\
    V0+TC       & 43 & 0.854          & 0.874          & 0.870          & 0.688          & 0.078          & 0.096          & $+0.037$\\
    V0+TA       & 43 & 0.860          & 0.882          & 0.872          & 0.656          & 0.081          & 0.101          & $-0.025$\\
    V0+TW       & 43 & 0.872          & 0.890          & 0.885          & 0.678          & 0.082          & 0.097          & $+0.026$\\
    V0+TWC      & 43 & 0.856          & 0.872          & 0.872          & 0.689          & 0.079          & 0.096          & $+0.064$\\
    V0+TWA      & 43 & 0.857          & 0.887          & 0.881          & 0.635          & 0.085          & 0.104          & $+0.000$\\
    V0+TWCA     & 43 & 0.869          & 0.876          & 0.871          & 0.704          & 0.075          & 0.093          & $-0.003$\\
    \textbf{V1} & 43 & \textbf{0.881} & \textbf{0.900} & \textbf{0.894} & \textbf{0.716} & \textbf{0.076} & \textbf{0.091} & $-0.026$\\
    \bottomrule
  \end{tabular*}
\end{table*}

\paragraph{V1 dominates on every metric.}
V1 is the only variant that simultaneously maximizes CCC, $r$,
Spearman $\rho$, and $R^2$ while minimizing MAE and RMSE. The two
elements that distinguish V1 from its standard-ablation predecessor
V0+TWCA are (a) per-dimension reference scores on each anchor
(beyond the holistic VC label) and (b) required per-dimension
numeric output. Because V1 uses the same model as every other cell,
the gain is a clean prompt-design effect: forcing the model to
commit to per-topic numbers, against per-topic anchors, sharpens
the holistic VC score even though that score is the only quantity
actually scored against ground truth. The bias column is the only
place V0 / V0+T / V0+TWCA edge V1 ($-0.026$ for V1 vs.\ near-zero
for several E1 cells), but their lower bias is not enough to
overcome larger correlation and scale-mismatch losses.

\paragraph{$R^2$ separates the variants more than CCC.}
$R^2$ against the identity line spreads variants more than CCC
does ($0.635$--$0.716$ vs.\ $0.852$--$0.881$): $R^2$ is more
sensitive to outlier residuals, so it discriminates between
variants whose CCC values are statistically tied. V0+TWCA
($R^2 = 0.704$) and V1 ($0.716$) are the clearest leaders here;
V0+A and V0+TWA sit at the bottom ($0.638$, $0.635$), suggesting
that anchor-only and weight+anchor cells produce a few large
residuals that the smoother CCC metric absorbs.

\subsection{Findings}
\label{app:vc-findings}

\begin{itemize}
  \item \textbf{Topics dominate, weighting helps, calibration and
    anchors do not stack.} The standard ablation peaks at V0+TW
    ($0.872$); adding C or A on top of V0+TW does not improve CCC.
  \item \textbf{Extended thinking interacts negatively with
    instruction density.} It helps V0 but hurts V0+TWCA and V0+TC,
    so adaptive thinking should be reserved for short prompts.
  \item \textbf{Per-dimension supervision is the only ingredient
    that breaks the $\sim\!0.87$ ceiling.} V1 wins on every
    agreement metric; the gain is a prompt-design effect, not a
    model-capacity effect.
\end{itemize}

\section{Topic Selection (Follow-on Study)}
\label{app:topicsel}

The 7 topics used inside the prompt (data density, visual encoding,
text legibility, schema, color/symbol, aesthetics, cognitive load)
were originally derived from a bottom-up coding of the same
user-comment corpus that produced the human VC ratings. Each topic
maps 1:1 to a scoring dimension in V1, so for every evaluation image
we also have a human-derived \emph{topic set} --- the subset of topics
that reviewers actually mentioned when describing that image. This
enables a second, complementary evaluation: rather than asking the
model to reproduce a continuous VC score, we ask whether it can
identify \emph{which topics} drive complexity for a given image.

\paragraph{Ground truth.} For each of the $n=43$ evaluation images,
we take the set of topics present in the compiled human comment pool
(the \texttt{Topics} column of \texttt{image\_compiled\_phrases.csv},
mapped to dimension keys). Human cardinality ranges 1--7 topics per
image (median 3); 14 images list $\geq 5$ topics.

\paragraph{Model task.} Variants that already expose the 7 topics to
the model (V0+T, V0+TW, V1) are extended to emit a top-3 topic
ranking. For V0+T and V0+TW the prompt is modified to request a
\texttt{top3\_topics} field (keys drawn from the canonical 7). V1
requires no prompt change: we derive top-3 directly from its
per-dimension numeric scores (highest-3 dims).

\paragraph{Metric.} Because human topic sets are unordered and
variable-sized, we report \textbf{set-based $F_1$} as the headline
($F_1 = 2PR/(P+R)$ with $P = |\text{pred} \cap \text{human}| / 3$ and
$R = |\text{pred} \cap \text{human}| / |\text{human}|$), plus
Jaccard/IoU, macro precision/recall (mean per image), and micro $F_1$
(pooled counts). Per-topic precision/recall/$F_1$ serve as
diagnostics.

\subsection{Aggregate results (n=43, opus-4.6)}

Table~\ref{tab:topicsel-agg} summarizes the top-3 topic-selection
performance for the three variants that already expose the 7-topic
taxonomy. Two patterns dominate the column.

\paragraph{Prompting for top-3 beats deriving from per-dimension
scores.} V0+T and V0+TW, which are asked explicitly to return a
top-3 topic list, both outperform the V1 runs by $\geq 9$ $F_1$
points (V0+TW $0.528$ vs.\ V1 $0.407$). V1 is strictly better at
the \emph{continuous} VC task (Sec.~\ref{app:vc-full}), but its
per-dimension scores --- trained to serve the holistic VC aggregate
--- do not discriminate the top-3 topics as sharply as a prompt
that asks for them directly. The V1-vs-E1 ordering of the
continuous task therefore reverses on the categorical task.

\paragraph{Qualitative weight guidance helps in aggregate.}
V0+TW exceeds V0+T by $+0.031$ $F_1$ ($0.528$ vs.\ $0.497$) on the
aggregate metric and is the leader on every column except Macro P /
Macro R. The ``high/medium/low'' bucketing biases selections toward
data\_density, visual\_encoding, schema, and cognitive\_load, which
are the four most-frequent human topics in the eval set. Whether
this aggregate gain is a free lunch is the question we take up in
the per-topic breakdown below.

\paragraph{$t=0$ and $t=1$-mean produce identical V1 results.}
The two V1 rows agree to four decimals on every metric, confirming
that the rank-of-three ordering is invariant to $t=1$ sampling
noise across the two V1 replicates. Sampling temperature is a
scoring-precision knob, not a ranking-structure knob, for this
task.

\begin{table}[h]
\centering
\caption{Topic-selection aggregate metrics for the top-3 task on
$n = 43$ non-anchor images, \texttt{claude-opus-4-6}. Each image has a
human topic set (1--7 topics, median 3) derived from the compiled user
comment pool. Predictions are a top-3 ranking over the canonical 7
topics. $F_1$ (headline, macro over images) is the mean over images
of $F_1 = 2PR/(P+R)$ with $P = |\text{pred} \cap \text{human}|/3$ and
$R = |\text{pred} \cap \text{human}|/|\text{human}|$;
IoU is the Jaccard index; Macro P/R are per-image means; Micro $F_1$
pools counts across images. V0+T and V0+TW emit top-3 directly via a
\texttt{top3\_topics} prompt field; V1 derives top-3 from its
per-dimension numeric scores (highest-3 dims). Bold = best per
column.}
\label{tab:topicsel-agg}
\scriptsize
\begin{tabular*}{\linewidth}{@{\extracolsep{\fill}}lccccc@{}}
\toprule
Variant & $F_1$ & IoU & Macro P & Macro R & Micro $F_1$ \\
\midrule
\textbf{V0+TW}      & \textbf{0.528} & \textbf{0.385} & 0.597 & 0.550 & \textbf{0.544} \\
V0+T                & 0.497 & 0.360 & 0.574 & 0.503 & 0.523 \\
V1 ($t=0$)          & 0.407 & 0.281 & 0.488 & 0.387 & 0.445 \\
V1 ($t=1$, mean)    & 0.407 & 0.281 & 0.488 & 0.387 & 0.445 \\
\bottomrule
\end{tabular*}
\end{table}

\subsection{Per-topic $F_1$}

The aggregate columns of Tab.~\ref{tab:topicsel-agg} hide a more
uneven picture across the 7 topics, shown in
Tab.~\ref{tab:topicsel-pertopic}.

\paragraph{Weighting trades per-topic fidelity for aggregate fit.}
V0+TW's qualitative weight guidance boosts \texttt{data\_density}
($0.765 \to 0.800$) and \texttt{cognitive\_load} ($0.566 \to 0.806$),
but actively \emph{suppresses} \texttt{text\_annotation}
($0.364 \to 0.087$, which is explicitly labeled ``low weight''),
and cuts \texttt{color\_symbol} from $0.538$ to $0.293$. V0+TW's
aggregate lead therefore comes from its alignment with the
topic-frequency distribution of the eval set, not from uniformly
better perception; in a deployment where every topic matters
equally, V0+T is the safer choice.

\paragraph{\texttt{aesthetic\_order} is unreachable across all
variants ($F_1 = 0$).} No variant ever ranks it in the top-3. It is
the rarest human topic ($6$ of $43$ images) and semantically
conflates with \texttt{data\_density} and \texttt{cognitive\_load}
in the model's representation; recovering it likely requires either
an explicit aesthetic-only sub-prompt or a different
ground-truth aggregation.

\paragraph{V0+T uniquely picks up \texttt{color\_symbol}.}
V0+T scores $F_1 = 0.538$ on color\_symbol, while V0+TW and V1
under-select it ($F_1 \approx 0.30$--$0.34$, high precision but low
recall). Color-driven complexity is easily absorbed into
\texttt{data\_density} or \texttt{visual\_encoding} by variants
that are not explicitly prompted to separate them.

\begin{table}[h]
\centering
\caption{Per-topic $F_1$ for the top-3 topic-selection task,
$n = 43$ images, \texttt{claude-opus-4-6}. Each cell is the $F_1$
computed over the single topic (treating prediction and ground truth
as binary: topic $\in$ set or not). Columns (abbreviated) from left to
right: \texttt{data\_density}, \texttt{visual\_encoding},
\texttt{text\_annotation}, \texttt{domain\_schema},
\texttt{color\_symbol}, \texttt{aesthetic\_order},
\texttt{cognitive\_load}. Bold per-column maxima.}
\label{tab:topicsel-pertopic}
\scriptsize
\setlength{\tabcolsep}{3pt}
\begin{tabular*}{\linewidth}{@{\extracolsep{\fill}}lccccccc@{}}
\toprule
Variant & dens. & vis.enc. & text & schema & color & aesth. & cog. \\
\midrule
V0+T           & 0.765 & 0.323 & 0.364 & 0.432 & \textbf{0.538} & 0.000 & 0.566 \\
\textbf{V0+TW} & \textbf{0.800} & 0.389 & 0.087 & 0.410 & 0.293 & 0.000 & \textbf{0.806} \\
V1 ($t=0$)     & 0.677 & 0.372 & 0.343 & \textbf{0.465} & 0.342 & 0.000 & 0.468 \\
V1 ($t=1$)     & 0.698 & 0.326 & 0.343 & \textbf{0.465} & 0.342 & 0.000 & 0.468 \\
\bottomrule
\end{tabular*}
\end{table}

\subsection{Per-vis-type \texorpdfstring{$F_1$}{F1} (n\,=\,63)}
\label{app:topicsel-pervistype}

To test whether prompt performance varies by vis-type, we expanded
the evaluation set to 63 images (7 per vis-type $\times$ 9 vis-types;
3 anchor images excluded) and re-evaluated V0+T and V0+TW.
Table~\ref{tab:topicsel-pervistype} shows precision, recall, and
$F_1$ per vis-type.

\textbf{Point plots are easiest; Glyph is hardest.}
Point achieves the highest $F_1$ for both variants (V0+T: 0.752,
V0+TW: 0.711), while Glyph is the lowest (0.401 / 0.458). Glyph
visualizations combine dense encoding with domain-specific spatial
arrangements, making their dominant topics harder to rank
unambiguously.

\textbf{V0+TW's weight guidance helps chart-centric types.}
V0+TW outperforms V0+T on Bar (+0.062), Cont.-ColorPatn (+0.071),
Glyph (+0.057), and Node-link (+0.036), all of which rely heavily on
the topics the weighting promotes (\texttt{data\_density},
\texttt{visual\_encoding}, \texttt{domain\_schema}).
Conversely, V0+T retains an advantage on Line (+0.058), Point
(+0.041), and Text (+0.052), where suppressing
\texttt{text\_annotation} and \texttt{color\_symbol} in V0+TW hurts.
Area and Grid are tied between variants.

\begin{table}[h]
\centering
\caption{Per-vis-type $F_1$, precision, and recall for the top-3
  topic-selection task, $n = 63$ images (7 per type, anchors excluded),
  \texttt{claude-opus-4-6}. \textbf{Bold}: higher $F_1$ per row.
  Bottom row: overall image-level macro-$F_1$.}
\label{tab:topicsel-pervistype}
\scriptsize
\setlength{\tabcolsep}{3pt}
\begin{tabular*}{\linewidth}{@{\extracolsep{\fill}}lc ccc ccc@{}}
\toprule
& & \multicolumn{3}{c}{V0+T} & \multicolumn{3}{c}{V0+TW} \\
\cmidrule(lr){3-5}\cmidrule(lr){6-8}
Vis-type & $n$ & $F_1$ & $P$ & $R$ & $F_1$ & $P$ & $R$ \\
\midrule
Area            & 7 & \textbf{0.463} & 0.571 & 0.466 & \textbf{0.463} & 0.571 & 0.466 \\
Bar             & 7 & 0.420 & 0.429 & 0.526 & \textbf{0.482} & 0.476 & 0.605 \\
Cont.-ColorPatn & 7 & 0.470 & 0.571 & 0.491 & \textbf{0.541} & 0.619 & 0.634 \\
Glyph           & 7 & 0.401 & 0.476 & 0.367 & \textbf{0.458} & 0.524 & 0.438 \\
Grid            & 7 & \textbf{0.533} & 0.619 & 0.514 & \textbf{0.533} & 0.619 & 0.514 \\
Line            & 7 & \textbf{0.580} & 0.619 & 0.571 & 0.522 & 0.571 & 0.500 \\
Node-link       & 7 & 0.535 & 0.524 & 0.676 & \textbf{0.571} & 0.571 & 0.705 \\
Point           & 7 & \textbf{0.752} & 0.857 & 0.705 & 0.711 & 0.810 & 0.669 \\
Text            & 7 & \textbf{0.468} & 0.571 & 0.421 & 0.416 & 0.524 & 0.356 \\
\midrule
Macro avg       & 63& 0.513 & & & \textbf{0.522} & & \\
\bottomrule
\end{tabular*}
\end{table}

\subsection{Findings}

\begin{itemize}
  \item \textbf{Categorical and continuous tasks favour different
    prompts.} V1 wins the continuous VC task; V0+TW wins the
    top-3 topic-selection task. Per-dimension numeric output
    sharpens scoring but does not sharpen ranking.
  \item \textbf{Weighting buys aggregate fit at the cost of
    per-topic fairness.} V0+TW leads on $F_1$ but suppresses the
    topics it labels ``low weight'' (text\_annotation, color\_symbol);
    pick V0+T when uniform per-topic recall matters more than
    aggregate $F_1$.
  \item \textbf{\texttt{aesthetic\_order} is a structural blind
    spot.} No variant recovers it; this is a limitation of the
    7-topic taxonomy and the model's lexical mapping, not of any
    individual prompt.
  \item \textbf{Vis-type moderates prompt choice.} On Glyph,
    Cont.-ColorPatn, Bar, and Node-link, V0+TW's weighting
    guidance lifts $F_1$ by 0.04--0.07; on Line, Point, and Text,
    the unweighted V0+T holds an equivalent advantage. No single
    variant dominates uniformly across vis-types.
\end{itemize}

% =========================================================================
% Appendix D: Structured Extraction --- Human-Phrase vs. Vision Grounding
% =========================================================================

\section{Structured Extraction: Human-Phrase vs.\ Vision Grounding}
\label{app:oar}

The experiments in Appendices~\ref{app:vc-prompts} and onwards evaluate
the model on the \emph{scalar} VC task. This section asks a
complementary question: when the model must produce structured scene
content --- objects, attributes, and relationships (O/A/R) --- how
close does vision-only grounding come to the structured content a
human-phrase baseline produces?

The extraction schema --- objects with regions, attributes bound to
objects, and relationships expressed as
$(\text{subject}, \text{predicate}, \text{object})$ triples ---
follows the \textbf{Visual Genome} dense-annotation formulation of
Krishna et al.~\cite{krishna2017vg}, which showed that objects /
attributes / relationships form a sufficient substrate for question
answering and grounded reasoning over natural images. We adapt that
formulation from natural scenes to data visualizations: objects are
chart components (bars, axes, legends, titles), attributes carry a
perceptual sentiment $\{+, -\}$ on the topic they instantiate, and
predicates encode complexity-relevant relations
(\texttt{increases\_effort}, \texttt{aids\_interpretation},
\texttt{obscures}, \texttt{simplifies}, \texttt{describes},
\texttt{structures}, \dots) rather than the spatial/action relations
(\texttt{riding}, \texttt{holding}) that dominate Visual Genome.
Canonicalization --- collapsing surface strings to a small vocabulary
of synsets and canonical predicates before any set-based comparison
--- is also adopted from the Visual Genome pipeline, where it was
necessary to make annotator-written triples comparable across 108k
images. The canonicalization dictionaries used here are defined in
\texttt{scripts/\_vc\_canon.py} and shared with the main pipeline
(\texttt{VCGenome\_FullPipeline.ipynb}).

\subsection{Three Conditions (n=43)}
\label{app:oar-conditions}

Using the same 43-image eval set (46 minus the 3 anchors from
Appendix~\ref{app:vc-gt}), we run three extraction conditions with
the same schema, the same 7-topic taxonomy, and the same decoding
settings ($t=0$, \texttt{claude-sonnet-4-6}, 1024 tokens):

\begin{table}[h]
\centering
\caption{Three O/A/R extraction conditions. B is the human-phrase
reference; V1 is pure vision; V2 adds three few-shot exemplars drawn
from B's own output on the anchor images.}
\label{tab:oar-conditions}
\footnotesize
\begin{tabular*}{\linewidth}{@{\extracolsep{\fill}}lll@{}}
\toprule
Cond. & Grounding signal & Input \\
\midrule
\textbf{B}  & topics + sentiment-tagged participant phrases & text only \\
\textbf{V1} & topics + image                           & vision \\
\textbf{V2} & topics + image + 3 anchors from B        & vision + exemplars \\
\bottomrule
\end{tabular*}
\end{table}

B serves as the human-phrase-grounded reference; V1 is pure vision;
V2 adds three in-context O/A/R exemplars taken from B's own output on
the three anchor images (so V2's anchors are guaranteed to be in the
same vocabulary and abstraction level as B). The anchor images
themselves are excluded from all metrics.

\subsection{Matching Procedure}
\label{app:oar-matching}

Each extraction is canonicalized with the pipeline's synset
dictionaries: objects are mapped to one of $\approx$40 synsets
(\texttt{mark.bar}, \texttt{structure.region}, \texttt{text.title},
\dots), predicates to one of $\approx$30 canonical relations, and
attributes are lowercased and snake-cased. We then compute set-based
precision / recall / $F_1$ / Jaccard per image and average over the
43 images. Three layers are reported: \emph{objects} (shared
synsets); \emph{attributes with sentiment} ($(\text{object\_synset},
\text{sentiment})$ pairs, wording-agnostic); and \emph{relationships}
in three granularities --- a \emph{strict} triple $(s, p, o)$, a
\emph{family} triple that rolls $p$ into one of $\approx$10 predicate
families (aids, hinders, obscures, simplifies,
increases\_complexity, increases\_effort, describes, structures,
\dots), and a \emph{loose} endpoint pair $\{s, o\}$.

\textbf{Role-level layer (cross-vocabulary alignment).} Preliminary
synset-level numbers revealed that B and the vision conditions
describe the same pictures with systematically different abstraction
levels: B names roles (\texttt{chart}, \texttt{data\_area},
\texttt{color\_encoding}, \texttt{annotations}) while V1/V2 name
instances (\texttt{horizontal\_bars}, \texttt{deceased\_bars},
\texttt{country\_labels}). To separate \emph{vocabulary mismatch}
from \emph{genuine disagreement}, we additionally project every
canonical synset to one of 8 chart-part \textbf{roles} ---
\texttt{chart}, \texttt{data\_area}, \texttt{encoding},
\texttt{title}, \texttt{axes}, \texttt{labels}, \texttt{legend},
\texttt{background} --- and recompute the three layers on the
projected keys. Role-level numbers reflect ``did the two extractions
agree on what chart region the claim is about,'' independent of
whether one says \texttt{horizontal\_bars} and the other says
\texttt{data\_area}.

\subsection{Results}
\label{app:oar-results}

We report mean per-image $F_1$ over the 43 evaluation images at
two levels of strictness. Tab.~\ref{tab:oar-strict} matches on
raw extracted keys (synsets and predicates), while
Tab.~\ref{tab:oar-role} first projects every synset to one of the
eight chart-part roles before matching, isolating conceptual
agreement from vocabulary choice.

\paragraph{Strict-key agreement.}
At the synset level, the two vision conditions (V1, V2) agree
strongly on objects ($F_1 = 0.68$) and on attribute sign
($F_1 = 0.59$), confirming that the model's perceptual lexicon for
chart components is internally consistent across two independently
prompted runs. Both vision conditions are far from the
human-phrase baseline B, however: B$\leftrightarrow$V object
$F_1$ is only $0.23$--$0.24$ and attribute-sign $F_1$ is
$0.11$. Relationships are essentially unaligned in every
pair: B$\leftrightarrow$V rel-family $F_1$ is $0.00$--$0.01$ and even
V1$\leftrightarrow$V2 rel-family is only $0.20$. A naive reading
would conclude that human-phrase summaries and vision-grounded
summaries describe almost disjoint scenes.

\begin{table}[h]
\centering
\caption{Mean $F_1$ across 43 non-anchor images, strict
synset-level keys. \emph{Attr (sign)} matches on
$(\text{object\_synset}, \text{sentiment})$;
\emph{Rel-family} rolls the predicate into $\approx$10
families; \emph{Rel-loose} ignores the predicate entirely.}
\label{tab:oar-strict}
\footnotesize
\setlength{\tabcolsep}{4pt}
\begin{tabular*}{\linewidth}{@{\extracolsep{\fill}}lcccc@{}}
\toprule
Pair      & Obj  & Attr (sign) & Rel-family & Rel-loose \\
\midrule
B vs V1   & 0.15 & 0.09        & 0.00       & 0.01 \\
B vs V2   & 0.15 & 0.09        & 0.00       & 0.01 \\
V1 vs V2  & 0.66 & 0.60        & 0.19       & 0.39 \\
\bottomrule
\end{tabular*}
\end{table}

\paragraph{Role-level agreement.}
Projecting synsets to eight chart-part roles
(\texttt{chart}, \texttt{data\_area}, \texttt{encoding},
\texttt{title}, \texttt{axes}, \texttt{labels}, \texttt{legend},
\texttt{background}) tells a different story
(Tab.~\ref{tab:oar-role}). The B$\leftrightarrow$V object $F_1$
roughly doubles --- $0.15 \to 0.25$ for B$\leftrightarrow$V1 and
$0.15 \to 0.28$ for B$\leftrightarrow$V2 (about $+67$--$+87\%$). The
two vision conditions also rise but by much less ($0.66 \to 0.72$),
because they already share an instance-level perceptual lexicon.
Most of the apparent B$\leftrightarrow$V gap on objects is
therefore \emph{vocabulary abstraction}, not disagreement about
what is in the chart: human commenters describe ``the chart'' or
``the labels,'' while the vision model names the specific instance
(``donut\_chart,'' ``slice\_labels''). Once both are collapsed to
the same eight roles they coincide much of the time.

\paragraph{Effect of phrase-grounded anchors.}
When B is grounded in the raw, sentiment-tagged participant phrases
(rather than a normalised curated set), the benefit of three
B-sourced anchors in V2 is small and not uniform. V2 narrowly edges
V1 on role-level objects ($0.25 \to 0.28$, $+0.03$) while
attribute-sign is tied at $0.19$; the corresponding
strict-synset deltas are within rounding ($0.00$ on both objects
and attribute-sign), and role-level loose relationships slightly
favour V1 ($0.06$ vs.\ $0.05$). Three exemplars are therefore
enough to nudge object vocabulary toward the human-phrase
distribution but do not shift attribute-sentiment or relationship
structure in any practically meaningful way. The implication is
that the B$\leftrightarrow$V gap is dominated by a generative
difference (what the human-phrase prompt elicits vs.\ what the
vision prompt elicits), which a handful of in-context exemplars
cannot close.

\begin{table}[h]
\centering
\caption{Mean $F_1$ across 43 non-anchor images, role-level keys:
every synset is projected to one of 8 chart-part roles before
matching. Bold = best B$\leftrightarrow$V entry per column.}
\label{tab:oar-role}
\footnotesize
\setlength{\tabcolsep}{3pt}
\begin{tabular*}{\linewidth}{@{\extracolsep{\fill}}lcccc@{}}
\toprule
Pair      & Obj-role      & Attr-role     & Rel-role-family & Rel-role-loose \\
          &               & (sign)        &                 &                \\
\midrule
B vs V1   & 0.25          & 0.19          & \textbf{0.03}   & \textbf{0.06} \\
B vs V2   & \textbf{0.28} & \textbf{0.19} & 0.02            & 0.05 \\
V1 vs V2  & 0.72          & 0.68          & 0.21            & 0.48 \\
\bottomrule
\end{tabular*}
\end{table}

\paragraph{Relationships remain the hardest layer.}
Unlike objects and attributes, the relationship layer barely
responds to role projection: B$\leftrightarrow$V rel-family $F_1$
stays at $0.02$--$0.03$ and rel-loose at $0.05$--$0.06$. Even the
two vision conditions agree on role-family relationships only at
$F_1 = 0.21$. Section~\ref{app:oar-example} shows that the
residual gap is conceptual rather than lexical.

\subsection{Example Scene Graphs}
\label{app:oar-example}

To illustrate the qualitative difference between the two grounding
strategies, Fig.~\ref{fig:oar-example} renders the canonicalized
scene graphs that B and V2 produced for the same image
(\texttt{SciVisJ.2926.1.png}, a 3D scientific flow visualization).
The two graphs are deliberately small
(B: 1/3/0, V2: 5/9/5 objects/attributes/relationships) to make the
vocabulary contrast easy to read.

\begin{figure*}[t]
  \centering
  \begin{minipage}[t]{0.45\linewidth}
    \centering
    \includegraphics[width=\linewidth]{figs/SciVisJ.2926.1.png}\\
    \small\textit{Source image}
  \end{minipage}%
  \hfill%
  \begin{minipage}[t]{0.40\linewidth}
    \centering
    \includegraphics[width=\linewidth]{figs/oar_SciVisJ.2926.1_B.pdf}\\
    \small\textit{Condition B (phrase-grounded)}
  \end{minipage}\\[8pt]
  \includegraphics[width=\linewidth]{figs/oar_SciVisJ.2926.1_V2.pdf}
  \caption{Canonicalized scene graphs for \texttt{SciVisJ.2926.1.png}.
    \textbf{Top left:} Source visualization image.
    \textbf{Top right:} Condition B (human-phrase grounded) uses generic
    chart-part roles and \emph{effectual}
    relationships (\texttt{fails\_to\_convey},
    \texttt{requires\_domain\_knowledge}). B's sole input phrase is
    \textit{``don't know what it presents.''~(+)}, so it cannot see the
    image and reaches for generic parts it uses for any
    chart. \textbf{Bottom:} Condition V2 (vision + B-sourced anchors)
    uses domain-specific scientific-visualization terms
    (\texttt{vortex\_structure}, \texttt{glyph\_spikes},
    \texttt{flow\_field}, \texttt{annotation\_label}) and
    \emph{structural} relationships (\texttt{contains},
    \texttt{overlaps\_with}, \texttt{describes},
    \texttt{requires\_expertise}). Node colors: green = object,
    purple = attribute (dark $=+$, light $=-$), red = relationship.}
  \label{fig:oar-example}
\end{figure*}

\subsection{Findings}
\label{app:oar-findings}

\begin{itemize}
  \item \textbf{Relationship divergence is conceptual, not lexical.}
    V1/V2 relationships are predominantly \emph{structural}
    (``bars stacked with bars,'' ``labels label rows,'' ``x-axis
    provides scale for bars''), while B relationships are
    predominantly \emph{effectual} (``color encoding adds visual
    complexity to data area,'' ``axis labels support readability of
    data area,'' ``data density increases reading effort of
    chart''); the contrast is visible in
    Fig.~\ref{fig:oar-example}. The human-phrase prompt pushes the
    model to reason about \emph{effects on perception}; the vision
    prompt pushes it to describe \emph{spatial arrangement}.
    Phrase-grounded anchors barely close this gap
    (rel-role-loose: V1 $0.12$ vs.\ V2 $0.13$), suggesting that
    in-context exemplars influence \emph{which entities} the model
    talks about much more than \emph{how} it relates them.
  \item \textbf{Headline takeaway for VC modelling.} Structured
    O/A/R extraction from images can recover the human baseline's
    \emph{what} and \emph{valence} only modestly once vocabulary is
    normalised (object-role $F_1 \approx 0.4$, attribute-sign role
    $F_1 \approx 0.25$ for V2), and it does \emph{not} recover the
    human baseline's \emph{causal framing} of the scene. Any
    downstream VC predictor that consumes O/A/R should therefore
    rely on objects and attribute-sentiment derived from vision, and
    treat the relationship layer either as human-phrase supervision
    or as a derived feature built on top of the canonical predicate
    families.
\end{itemize}

% =========================================================================
% Appendix E: Full 510-Image Evaluation
% =========================================================================

\section{VC Prediction on the Full Dataset: V0+TW vs.\ V1}
\label{app:fullset}

The ablations in Appendices~\ref{app:vc-prompts}--\ref{app:topicsel}
and the extraction study in Appendix~\ref{app:oar} all operate on the
46-image pilot set ($n = 43$ after anchor exclusion). That set was
deliberately constructed to maximise $|\Delta\mathrm{VC}|$ coverage
rather than to be representative of the broader image pool --- its wider
ground-truth spread inflates CCC relative to an unselected sample.
This section extends the VC prediction evaluation to the full curated
image corpus to assess how the two best-performing prompts generalise:
\textbf{V0+TW} (the strongest deterministic E1 cell,
Sec.~\ref{app:vc-ccc}) and \textbf{V1} (the upper envelope of the
ablation).

\subsection{Setup}
\label{app:fullset-setup}

\paragraph{Image set.}
We use human VC ratings from the same pair-comparison study described in
Sec.~\ref{app:vc-gt}, extended to the full curated image pool of
approximately 520 images with complete annotation. Retaining the nine
vis-types with sufficient representation (Bar, Line, Area, Point,
Node-link, Glyph, Grid, Cont.-ColorPatn, Text) and dropping four
rare categories (Schematic, Bar\,\&\,Point, Table, Area\,\&\,Text) yields
\textbf{510 images}. Following the same leakage-prevention convention
applied in Sec.~\ref{app:vc-gt}, the three anchor exemplars are excluded
from both prompts' metrics, yielding $n = 507$ test images. V1 requires
this exclusion because it receives the anchors as in-context examples;
we apply it to V0+TW as well so both conditions are evaluated on an
identical image set.
Full-set ground-truth statistics: mean $= 0.598$, SD $= 0.148$ --- a
noticeably narrower spread than the pilot (pilot SD $\approx 0.244$),
consistent with the pilot's deliberate difficulty stratification.

\paragraph{Runs.}
Both prompts run on \texttt{claude-opus-4-6} at $t = 0$ (deterministic).
V1 was originally scored at the API default sampling temperature ($t = 1$)
on the unfiltered 520-image pool; to eliminate a temperature confound
in the prompt comparison, we re-ran V1 at $t = 0$ on the same 510-image
set used for V0+TW. The temperature sensitivity analysis confirming that
this re-run does not materially change V1's output is reported in
Sec.~\ref{app:fullset-temp}. One deterministic run per prompt is
available (no double-replicate averaging as in
Sec.~\ref{app:vc-conditions}); all figures below are single-run point
estimates.

\subsection{Overall Results}
\label{app:fullset-overall}

Table~\ref{tab:fullset-overall} reports the full agreement metrics for
both prompts. The most important observation is the headline drop
relative to the pilot: V0+TW falls from CCC $= 0.872$ to $0.786$ and
V1 from $0.881$ to $0.780$. The decline is expected --- the pilot's
easy-discrimination pairs inflate the CCC denominator, and the full
set re-introduces many mid-range images where small residuals
accumulate. The drop is proportionally similar for both prompts,
suggesting it reflects the evaluation-set shift rather than a
prompt-specific failure.

\begin{table*}[t]
  \centering
  \caption{Full agreement metrics for V0+TW and V1 on $n = 507$
    images (the three anchor exemplars are excluded from both conditions
    for a consistent comparison), \texttt{claude-opus-4-6}, $t = 0$.
    CCC is Lin's Concordance Correlation Coefficient (headline; jointly
    penalises low correlation, mean bias, and scale mismatch). $r$ is
    Pearson; $\rho$ is Spearman; $R^2$ is the coefficient of
    determination against the identity line (can be negative if
    predictions are biased off the diagonal). MAE and RMSE are in the
    raw $[0,1]$ VC scale; Bias is the signed mean error
    $\overline{\hat{y}-y}$. All Pearson $p < 10^{-12}$.}
  \label{tab:fullset-overall}
  \small
  \begin{tabular*}{\textwidth}{@{\extracolsep{\fill}}lcccccccc@{}}
    \toprule
    Prompt & $t$ & $n$ & CCC & $r$ & $\rho$ & $R^2$ & MAE & RMSE & Bias \\
    \midrule
    V0+TW & 0 & 507 & 0.784 & 0.835 & 0.833 & 0.370 & 0.093 & 0.116 & $-0.030$ \\
    V1    & 0 & 507 & 0.777 & 0.854 & 0.851 & 0.388 & 0.093 & 0.115 & $-0.063$ \\
    \bottomrule
  \end{tabular*}
\end{table*}

\paragraph{The two prompts are indistinguishable at the headline level.}
On the identical 507-image test set, V0+TW and V1 differ by
$\Delta\mathrm{CCC} = -0.007$ and $\Delta\mathrm{MAE} = 0.000$.
V1 gains approximately $0.02$ on Pearson $r$, but this improvement is
offset by a larger negative bias: V1 under-predicts by $-0.063$ vs.\
V0+TW's $-0.030$. Because CCC jointly penalises mean shift alongside
correlation loss, the net CCC difference is negligible.

\paragraph{Prediction spread.}
Both prompts over-spread relative to the human ground truth (GT SD
$= 0.148$): V0+TW's predicted SD is $0.201$ and V1's is $0.183$. V1's
tighter spread partly explains its higher Pearson $r$ (a scale-invariant
metric) while its larger negative bias neutralises any CCC advantage.

\paragraph{Note on $R^2$ conventions.}
Throughout this appendix, $R^2$ is \texttt{sklearn.metrics.r2\_score}
evaluated against the identity line $\hat{y} = y$, which can go negative
when predictions are biased off the diagonal. An alternative convention
reports the squared Pearson correlation $r^2$ from an unconstrained
linear regression of predictions on ground truth, which is bounded in
$[0,1]$ and does not penalise mean shift or scale mismatch. For V1 at
$t = 0$ on the 507-image test set, the two conventions give $r^2 = 0.730$
vs.\ identity-line $R^2 = 0.388$; the gap of $0.342$ is entirely
attributable to bias ($-0.063$) and scale mismatch
($\hat\sigma / \sigma_y = 1.24$).

\subsection{Per-Vis-Type Breakdown}
\label{app:fullset-vistype}

Table~\ref{tab:fullset-vistype} breaks the 510-image results down by
vis-type. No prompt dominates uniformly across all nine types; the
outcomes are split.

\begin{table*}[t]
  \centering
  \caption{Per-vis-type agreement metrics on the 507-image test set
    (anchors excluded), \texttt{claude-opus-4-6}, $t = 0$. CCC is the
    headline; $R^2$ is the identity-line score (can be negative);
    MAE is in raw $[0,1]$ VC units. $\Delta$CCC $=$ CCC(V1) $-$
    CCC(V0+TW); positive entries favour V1.
    Types are ordered by sample size.}
  \label{tab:fullset-vistype}
  \small
  \begin{tabular*}{\textwidth}{@{\extracolsep{\fill}}lrccccccc@{}}
    \toprule
    Vis Type & $n$ & CCC (V0+TW) & CCC (V1) & $\Delta$CCC
             & $R^2$ (V0+TW) & $R^2$ (V1) & MAE (V0+TW) & MAE (V1) \\
    \midrule
    Node-link       &  66 & 0.797 & 0.818 & $+0.021$          & $\phantom{-}0.434$ & $\phantom{-}0.555$ & 0.096 & 0.080 \\
    Area            &  65 & 0.833 & 0.825 & $-0.008$          & $\phantom{-}0.527$ & $\phantom{-}0.519$ & 0.080 & 0.083 \\
    Grid            &  64 & 0.738 & 0.731 & $-0.007$          & $\phantom{-}0.275$ & $\phantom{-}0.272$ & 0.097 & 0.099 \\
    Glyph           &  64 & 0.729 & 0.719 & $-0.010$          & $\phantom{-}0.062$ & $\phantom{-}0.102$ & 0.102 & 0.101 \\
    Point           &  57 & 0.761 & 0.671 & $-0.090$          & $\phantom{-}0.299$ & $\phantom{-}0.011$ & 0.090 & 0.111 \\
    Bar             &  50 & 0.823 & 0.785 & $-0.038$          & $\phantom{-}0.514$ & $\phantom{-}0.445$ & 0.092 & 0.098 \\
    Text            &  51 & 0.781 & 0.761 & $-0.020$          & $\phantom{-}0.498$ & $\phantom{-}0.404$ & 0.070 & 0.080 \\
    Line            &  48 & 0.732 & 0.745 & $+0.013$          & $\phantom{-}0.164$ & $\phantom{-}0.292$ & 0.103 & 0.090 \\
    Cont.-ColorPatn &  42 & 0.562 & 0.614 & $+0.052$          & $-1.223$           & $-0.540$           & 0.116 & 0.098 \\
    \bottomrule
  \end{tabular*}
\end{table*}

\paragraph{Best-performing types (Bar, Area).}
Both prompts achieve CCC $\geq 0.826$ on Bar and Area --- the two most
common statistical chart categories, where the model's prior apparently
aligns well with human VC ratings and the complexity gradient is easy to
operationalise.

\paragraph{V1 wins on Node-link and Cont.-ColorPatn.}
V1's gains are largest on Cont.-ColorPatn ($+0.052$ CCC) and Node-link
($+0.021$). For Cont.-ColorPatn the ground-truth spread is unusually
narrow (GT SD $= 0.105$), so any over-spread prediction falls
systematically off the identity line; V1's better-calibrated predictions
(pred SD V0+TW $\approx 0.203$ vs.\ V1 $\approx 0.167$) convert a
catastrophic $R^2 = -1.22$ into a merely poor $R^2 = -0.54$ while
lifting CCC by $0.052$. Despite this improvement, both prompts produce
$R^2 < 0$ on this type, making Cont.-ColorPatn a structural weakness.

\paragraph{V0+TW wins on Point and Bar.}
The largest V1 regression is on Point charts, where V1's CCC drops from
$0.754$ to $0.669$ ($\Delta = -0.085$) and $R^2$ collapses from
$0.281$ to $0.014$. The calibration ranges and anchor exemplars in V1
appear to push Point predictions away from the true mid-range cluster
for this type. V0+TW also holds a clear lead on Bar ($\Delta = -0.037$),
Text ($\Delta = -0.020$), and Area ($\Delta = -0.008$).

\subsection{Effect of Temperature}
\label{app:fullset-temp}

To verify that temperature is not a confound in the V0+TW vs.\ V1
comparison, we compare the original V1 run at $t = 1$ (unfiltered
520-image pool) with the V1 re-run at $t = 0$ (510-image set,
$n = 507$ after anchor exclusion) on the overlapping image subset. Across all agreement statistics (CCC, Pearson $r$,
$R^2$, MAE), the shift from $t = 1$ to $t = 0$ is $\leq 0.005$ ---
well within single-image sampling noise. Temperature setting is
therefore not a material confound in the V0+TW vs.\ V1 comparison.

\subsection{Findings}
\label{app:fullset-findings}

\begin{itemize}
  \item \textbf{Full-set CCC is lower than pilot CCC for both prompts}
    (V0+TW: $0.786$ vs.\ $0.872$; V1: $0.780$ vs.\ $0.881$), primarily
    because the pilot's difficulty-stratified sampling produces a wider
    ground-truth spread that flatters CCC.
  \item \textbf{V0+TW and V1 are indistinguishable at the headline
    level} ($\Delta\mathrm{CCC} = -0.007$, $\Delta\mathrm{MAE} = 0.000$
    on the identical $n = 507$ test set) with temperature and model held
    constant. The extra V1 machinery (Calibration + Anchors +
    per-dimension JSON output) does not yield a meaningful accuracy
    improvement on a representative image pool.
  \item \textbf{Per-type outcomes are split.} V1 is better for Node-link
    and Cont.-ColorPatn; V0+TW is better for Point and Bar. Applications
    sensitive to Point-chart accuracy should prefer V0+TW; those
    targeting Node-link or continuous colour maps may benefit from V1.
  \item \textbf{Cont.-ColorPatn is a structural weakness for both
    prompts} ($R^2 < 0$). Improving performance on this type likely
    requires either type-specific calibration guidance or a broader
    anchor set covering extreme values in the category.
  \item \textbf{The practical recommendation from
    Sec.~\ref{app:vc-findings} holds on the full set:} V0+TW is the
    preferred prompt --- it matches V1 on overall agreement with fewer
    instructions, no per-dimension output overhead, and better
    performance on the high-frequency Point and Bar types.
\end{itemize}

% ======== Section: Prompt Listings =======================================

\section{Prompt Listings}
\label{app:prompts}

All variants share a common structure: a \emph{system prompt} (delivered
once per API call) followed by one or more \emph{user turns} (one per
image). The system prompt is built by starting from the \textbf{V0}
baseline and inserting the four modular component blocks (\textbf{T},
\textbf{W}, \textbf{C}, \textbf{A}) described in
Sec.~\ref{app:vc-components}. Each block below is the exact text used
in the deployed prompt.

% ─── V0 baseline ──────────────────────────────────────────────────────────
\subsection{V0 Baseline (Complete System Prompt)}
\label{app:prompts-v0}

The following is the complete system prompt for \textbf{V0}. Every other
variant starts from this text and inserts one or more component blocks
immediately before the \texttt{\#\#~Output~Format} section.

\begin{lstlisting}
You are a visual complexity (VC) scoring expert for data visualizations.

**Visual Complexity** measures the amount of detail, intricacy, and perceptual/cognitive demand of a visualization image. It is scored on a 0-1 scale where 0 = minimal complexity and 1 = maximum complexity.

You will receive a single visualization image with NO accompanying text or comments. Score it purely from what you see.

## Output Format
Return ONLY valid JSON (no markdown fences, no explanation outside JSON):
{
  "vc_score": <float 0-1>,
  "explanation": "<2-3 sentence justification>"
}
\end{lstlisting}

\noindent The user turn for all non-anchored variants is a single image
block accompanied by the instruction ``Score the visual complexity of
this visualization image.''

% ─── Component blocks ─────────────────────────────────────────────────────
\subsection{Component Blocks}
\label{app:prompts-components}

Each block is inserted into the system prompt immediately before the
\texttt{\#\#~Output~Format} section, in the order T\,$\rightarrow$\,W\,$\rightarrow$\,C.
The \textbf{A} component does not add text to the system prompt; it
extends the message turn sequence instead (see
Sec.~\ref{app:prompts-anchors}).

\paragraph{Block T --- Topics / Dimensions.}
\textit{Used in: V0+T, V0+TC, V0+TA, V0+TW, V0+TWC, V0+TWA, V0+TWCA, V1.}

\begin{lstlisting}
## Topics to Consider When Scoring

When assessing visual complexity, consider the following 7 topics that commonly drive complexity perception. You do not need to score each topic individually -- use them as mental checkpoints to arrive at a holistic vc_score.

1. **Data Density / Image Clutter** -- The perceived amount, richness, or depth of data content. Considers information volume, element quantity, and visual clutter/overlap.
2. **Visual Encoding Clarity** -- The variety, type, and complexity of graphical forms (shapes, lines, marks) and how spatial layout, scale, and encoding interpretability contribute to complexity.
3. **Semantics / Text Legibility** -- The quantity and density of text elements (titles, axis labels, legends, captions, annotations, in-chart labels).
4. **Schema** -- Whether specialized domain knowledge is needed, including dimensionality (2D/3D), structural complexity, and abstraction level.
5. **Color, Symbol, and Texture Details** -- Range, variety, and arrangement of colors, plus use of symbols, textures, and non-color graphical markers.
6. **Aesthetics Uncertainty** -- How visually cluttered, dense, or disordered the layout appears. Higher = more cluttered/overwhelming. A clean minimal layout scores low; a crowded layout with overlapping elements scores high.
7. **Immediacy / Cognitive Load** -- Overall ease or difficulty of interpreting the visualization. Considers interpretive difficulty, semantic clarity, and processing time/effort.
\end{lstlisting}

\paragraph{Block W --- Weighted Scoring.}
Appended immediately after Block~T (requires T to be present).
\textit{Used in: V0+TW, V0+TWC, V0+TWA, V0+TWCA, V1.}

\begin{lstlisting}
## Weighting of Topics for the Overall vc_score
The overall vc_score is a weighted holistic judgment. Weight the 7 topics as follows:
- **High weight**: Data Density / Image Clutter, Visual Encoding Clarity, Schema, Immediacy / Cognitive Load (these drive VC the most).
- **Medium weight**: Color / Symbol / Texture Details, Aesthetics Uncertainty.
- **Low weight**: Semantics / Text Legibility (text volume alone is a weak VC signal).
A low score on Semantics / Text Legibility alone should NOT substantially pull down the overall vc_score.
\end{lstlisting}

\noindent When W is present, the \texttt{"vc\_score"} field description
in \texttt{\#\#~Output~Format} is also changed to
\texttt{<float~0-1,~weighted~holistic~judgment>}.

\paragraph{Block C --- Calibration.}
Appended after Block~T (and W if present), before
\texttt{\#\#~Output~Format}.
\textit{Used in: V0+C, V0+TC, V0+TWC, V0+TWCA, V1.}

\begin{lstlisting}
## Calibration Guidance
Use the full 0-1 range. Do NOT cluster scores conservatively toward the center.
- A plain single bar chart or pie chart with < 5 elements -> vc_score approx 0.25-0.40
- A standard scatterplot, grouped bar chart, or line chart -> vc_score approx 0.45-0.60
- A multi-panel dashboard, dense network, or multi-encoding composite -> vc_score approx 0.70-0.85
- An extremely dense multi-view system with many coordinated panels -> vc_score approx 0.85-0.95
IMPORTANT: If you are unsure between two scores, always choose the higher one. Most scorers underestimate -- push scores up.
\end{lstlisting}

\paragraph{Block A --- Anchor Few-Shot Turns.}
\label{app:prompts-anchors}
The A component does not modify the system prompt. Instead, it prepends
three user$\to$assistant turn pairs to the message list before the
target-image query. The three anchors span the scale: low (0.22), mid
(0.54), high (0.95).
\textit{Used in: V0+A, V0+TA, V0+TWA, V0+TWCA, V1.}

\begin{lstlisting}
[USER]   [IMAGE: VisC.503.6.png]
         Score the visual complexity of this visualization image.
[ASST]   {"vc_score": 0.22, "explanation": "Anchor image scored at 0.22."}

[USER]   [IMAGE: InfoVisJ.619.17.png]
         Score the visual complexity of this visualization image.
[ASST]   {"vc_score": 0.54, "explanation": "Anchor image scored at 0.54."}

[USER]   [IMAGE: InfoVisJ.1149.6(1).png]
         Score the visual complexity of this visualization image.
[ASST]   {"vc_score": 0.95, "explanation": "Anchor image scored at 0.95."}

[USER]   [IMAGE: target visualization]   <- actual query
         Score the visual complexity of this visualization image.
\end{lstlisting}

\noindent Anchor images are excluded from all reported metrics to
prevent leakage.

% ─── V1 extended ──────────────────────────────────────────────────────────
\subsection{V1: Extended Variant}
\label{app:prompts-v1}

\textbf{V1} uses the same system prompt as V0+TWCA (Base\,+\,T\,+\,W\,+\,C)
with two modifications.

\textbf{Modification 1 --- Per-dimension anchor scores.} Each anchor
assistant turn carries full dimension scores instead of only
\texttt{vc\_score}:

\begin{lstlisting}
[ASST]   {
           "data_density": 0.15, "visual_encoding": 0.20,
           "text_annotation": 0.10, "domain_schema": 0.15,
           "color_symbol": 0.10, "aesthetic_order": 0.20,
           "cognitive_load": 0.20, "vc_score": 0.22,
           "data_density_explanation": "...",
           ...
           "summary": "..."
         }
\end{lstlisting}

\textbf{Modification 2 --- Extended output schema.} The
\texttt{\#\#~Output~Format} block is replaced to require numeric
per-dimension output, and the section header becomes
\texttt{\#\#~Scoring~Dimensions~(each~0--1)}:

\begin{lstlisting}
## Output Format
Return ONLY valid JSON (no markdown fences, no explanation outside JSON):
{
  "data_density": <float 0-1>,
  "visual_encoding": <float 0-1>,
  "text_annotation": <float 0-1>,
  "domain_schema": <float 0-1>,
  "color_symbol": <float 0-1>,
  "aesthetic_order": <float 0-1>,
  "cognitive_load": <float 0-1>,
  "vc_score": <float 0-1>,
  "data_density_explanation": "<1 sentence>",
  "visual_encoding_explanation": "<1 sentence>",
  "text_annotation_explanation": "<1 sentence>",
  "domain_schema_explanation": "<1 sentence>",
  "color_symbol_explanation": "<1 sentence>",
  "aesthetic_order_explanation": "<1 sentence>",
  "cognitive_load_explanation": "<1 sentence>",
  "summary": "<2-3 sentence overall description>"
}
\end{lstlisting}

% ─── Topic matching ───────────────────────────────────────────────────────
\subsection{Topic Matching (Top-3 Selection)}
\label{app:prompts-topic}

The topic-matching study (Sec.~\ref{app:topicsel}) extends the V0+T and
V0+TW VC-scoring prompts with one additional output field,
\texttt{top3\_topics}, and a \textbf{Topic-Key Mapping} block. The system
prompt is otherwise identical to V0+T or V0+TW respectively; only the
\texttt{\#\#~Output~Format} section and one new block differ.

\paragraph{Block: Topic-Key Mapping.}
Inserted immediately before \texttt{\#\#~Output~Format}:

\begin{lstlisting}
## Topic-Key Mapping
Each topic corresponds to a short key:
- "Data Density / Image Clutter"     -> data_density
- "Visual Encoding Clarity"          -> visual_encoding
- "Semantics / Text Legibility"      -> text_annotation
- "Schema"                           -> domain_schema
- "Color, Symbol, and Texture Details" -> color_symbol
- "Aesthetics Uncertainty"           -> aesthetic_order
- "Immediacy / Cognitive Load"       -> cognitive_load
\end{lstlisting}

\paragraph{Modified Output Format (V0+T top-3 variant).}
Replaces the standard \texttt{\#\#~Output~Format} block:

\begin{lstlisting}
## Output Format
Return ONLY valid JSON (no markdown fences, no explanation outside JSON):
{
  "vc_score": <float 0-1>,
  "top3_topics": ["<key>", "<key>", "<key>"],
  "explanation": "<2-3 sentence justification referencing the topics that most influenced the score>"
}

Where `top3_topics` is an ordered list (most -> least influential) of exactly three distinct keys drawn from: data_density, visual_encoding, text_annotation, domain_schema, color_symbol, aesthetic_order, cognitive_load.
\end{lstlisting}

\noindent The V0+TW top-3 variant is identical except the \texttt{"vc\_score"} field description reads \texttt{<float~0-1,~weighted~holistic~judgment>} (consistent with Block~W). For \textbf{V1}, no prompt modification is needed: the top-3 topics are derived programmatically by taking the three highest per-dimension scores from the V1 JSON output.

% ─── OAR extraction ───────────────────────────────────────────────────────
\subsection{OAR Extraction}
\label{app:prompts-oar}

The structured O/A/R extraction study (Sec.~\ref{app:oar}) uses three
conditions that share a common \textbf{Topics block} and
\textbf{Output Schema block}; only the preamble paragraph differs per
condition.

\paragraph{Shared block: Topics (7-topic taxonomy).}
Appended in all three system prompts before the output schema.

\begin{lstlisting}
## Topics (the 7-topic taxonomy)

Use these 7 topics when labeling attributes and relationships. Each attribute and relationship must be tagged with exactly one topic.

1. **Data Density / Image Clutter** -- The perceived amount, richness, or depth of data content. Considers information volume, element quantity, and visual clutter/overlap.
2. **Visual Encoding Clarity** -- The variety, type, and complexity of graphical forms (shapes, lines, marks) and how spatial layout, scale, and encoding interpretability contribute to complexity.
3. **Semantics / Text Legibility** -- The quantity and density of text elements (titles, axis labels, legends, captions, annotations, in-chart labels).
4. **Schema** -- Whether specialized domain knowledge is needed, including dimensionality (2D/3D), structural complexity, and abstraction level.
5. **Color, Symbol, and Texture Details** -- Range, variety, and arrangement of colors, plus use of symbols, textures, and non-color graphical markers.
6. **Aesthetics Uncertainty** -- How visually cluttered, dense, or disordered the layout appears. Higher = more cluttered/overwhelming.
7. **Immediacy / Cognitive Load** -- Overall ease or difficulty of interpreting the visualization. Considers interpretive difficulty, semantic clarity, and processing time/effort.
\end{lstlisting}

\paragraph{Shared block: Output Schema.}
Identical across all three conditions. Note the field is \texttt{"topic"}
(not \texttt{"subtopic"}).

\begin{lstlisting}
## Output Schema

Return ONLY valid JSON (no markdown fences, no prose outside JSON):
{
  "objects": [
    {"id": 1, "name": "object_name", "region": "data_area|axes|legend|title|annotation|overall"}
  ],
  "attributes": [
    {"object_id": 1, "attr": "short_snake_case_phrase", "sentiment": "+" or "-", "topic": "<one of the 7 topics>"}
  ],
  "relationships": [
    {"subj": 1, "pred": "snake_case_predicate", "obj": 2, "sentiment": "+" or "-", "topic": "<one of the 7 topics>"}
  ]
}

Guidelines:
- Extract elements proportional to the richness of the input. Typical ranges: 2-6 objects, 3-8 attributes, 1-5 relationships; sparse inputs may produce fewer and rich inputs may exceed these.
- Each list must contain at least one entry.
- Object names: single lowercase word or snake_case.
- Attribute text: snake_case, max 4 words.
- Sentiment: '+' increases perceived complexity, '-' decreases it.
- Topic must be exactly one of the seven topic titles listed above.
\end{lstlisting}

\paragraph{Condition B --- Phrase-Grounded (no image).}
Input: sentiment-tagged complexity phrases from participant comments.

\begin{lstlisting}
You are a visual complexity annotation expert. You will receive, for a single data-visualization image:
  (a) a list of original complexity phrases extracted from real participant comments about the image, each tagged with a sentiment marker (+) (increases perceived complexity) or (-) (decreases perceived complexity).

Extract a scene graph (objects, attributes, relationships) that is **strictly grounded in the phrases**. Do not invent elements the phrases do not mention or imply. Use the (+) / (-) markers to inform the `sentiment` field of attributes and relationships.

[TOPICS BLOCK]

[OUTPUT SCHEMA BLOCK]
\end{lstlisting}

\noindent The user turn supplies the image filename and the
semicolon-separated phrase list with per-phrase sentiment tags.

\paragraph{Condition V1 --- Vision-Only (image, no phrases).}
Input: image only (no phrases).

\begin{lstlisting}
You are a visual complexity annotation expert. You will receive, for a single data-visualization image:
  (a) the rendered image.

Extract a scene graph (objects, attributes, relationships) that is **strictly grounded in what is visible in the image**. Do not invent elements that are not visible.

[TOPICS BLOCK]

[OUTPUT SCHEMA BLOCK]
\end{lstlisting}

\paragraph{Condition V2 --- Vision + Anchor Exemplars.}
\texttt{SYSTEM\_V2} appends one additional block to the V1 system prompt:

\begin{lstlisting}
## Reference Anchor Examples

Below are three reference images with expert-produced scene graphs. Use them to calibrate your extraction style (granularity, object naming, predicate style, topic usage). Do NOT copy their content; produce an extraction for the target image shown last.
\end{lstlisting}

\noindent The three anchor user$\to$assistant turns are taken from
Condition~B's output on the same three anchor images, keeping the
exemplar style consistent with the text-grounded reference. Anchor images
are excluded from all evaluation comparisons.

% ======== BibTeX entry for bibliography file ==============================
% @article{lin1989ccc,
%   author  = {Lin, Lawrence I-Kuei},
%   title   = {A Concordance Correlation Coefficient to Evaluate Reproducibility},
%   journal = {Biometrics},
%   volume  = {45},
%   number  = {1},
%   pages   = {255--268},
%   year    = {1989},
%   doi     = {10.2307/2532051},
%   publisher = {International Biometric Society},
% }
%
% @article{krishna2017vg,
%   author    = {Krishna, Ranjay and Zhu, Yuke and Groth, Oliver and
%                Johnson, Justin and Hata, Kenji and Kravitz, Joshua and
%                Chen, Stephanie and Kalantidis, Yannis and Li, Li-Jia and
%                Shamma, David A. and Bernstein, Michael S. and
%                Fei-Fei, Li},
%   title     = {Visual Genome: Connecting Language and Vision Using
%                Crowdsourced Dense Image Annotations},
%   journal   = {International Journal of Computer Vision},
%   volume    = {123},
%   number    = {1},
%   pages     = {32--73},
%   year      = {2017},
%   doi       = {10.1007/s11263-016-0981-7},
% }
